% *** ก่อนส่ง: ลบข้อความตัวอย่าง (\ldots) ทั้งหมด แล้วเขียนเนื้อหาของตนเอง ***
% =============================================================
% บทที่ 1 - บทนำ
% แก้ไขเนื้อหาในบทนี้ได้ที่ไฟล์ chapter1.tex
% =============================================================

% -----------------------------------------------------------
\section{ความเป็นมาและความสำคัญของปัญหา}
% -----------------------------------------------------------

% เขียนอธิบายที่มาของปัญหา ความสำคัญ และความจำเป็นในการพัฒนา
% ควรเริ่มจากปัญหาในภาพกว้าง แล้วค่อยๆ เจาะจงลงมา

ในปัจจุบัน ระบบ\ldots{} มีบทบาทสำคัญในด้าน\ldots{} อย่างไรก็ตาม ยังคงมีปัญหาเกี่ยวกับ\ldots{} ซึ่งส่งผลให้\ldots{} ดังนั้น จึงมีความจำเป็นต้องพัฒนา\ldots{} เพื่อแก้ไขปัญหาดังกล่าว

% -----------------------------------------------------------
\section{วัตถุประสงค์}
% -----------------------------------------------------------

% ระบุวัตถุประสงค์ของโครงงานเป็นข้อๆ
% ใช้คำกริยาที่ชัดเจน เช่น "เพื่อพัฒนา" "เพื่อศึกษา" "เพื่อประเมิน"

\begin{enumerate}[label=\arabic*)]
    \item เพื่อศึกษาและวิเคราะห์ปัญหา\ldots{}
    \item เพื่อออกแบบและพัฒนาระบบ\ldots{}
    \item เพื่อทดสอบและประเมินประสิทธิภาพ\ldots{}
\end{enumerate}


% % -----------------------------------------------------------
% \section{เป้าหมายและขอบเขต}
% % -----------------------------------------------------------

% \subsection{ขอบเขตด้านเนื้อหา}

% ระบบที่พัฒนาครอบคลุมการทำงานในส่วนของ\ldots{} โดยมุ่งเน้นที่\ldots{}

% \subsection{ขอบเขตด้านเวลาและทรัพยากร}

% โครงงานนี้ดำเนินการในระยะเวลา\ldots{} โดยใช้ทรัพยากร\ldots{}

% \subsection{ขอบเขตด้านเทคนิคหรือเครื่องมือ}

% พัฒนาด้วยภาษา\ldots{} บนแพลตฟอร์ม\ldots{} ร่วมกับฐานข้อมูล\ldots{}

% -----------------------------------------------------------
\section{ขอบเขตและข้อจำกัดของระบบ}
% -----------------------------------------------------------

\subsection{ขอบเขต}

ให้ภาพรวมของระบบ ระบุทั้ง Functional requirements และ Non-functional requirements ให้ครบถ้วน ในกรณีที่โครงงานเป็นแบบวิจัย ให้ระบุถึงกระบวนการทดลอง การเก็บช้อมูลหรือข้อมูลที่จะนำมาใช้ ภาพแวดล้อมในการทดสอบทดลอง และวิธีการวัดประเมินผลการทดลองแทน

\subsection{ข้อจำกัดของโครงงาน}
ระบุข้อจัดกัดหรือ Scope ของผลลัพธ์ที่จะได้จากโครงงานเช่น 
\begin{enumerate}[label=\arabic*)]
    \item เว็บแอปพลิเคชันสามารถแสดงผลได้ 2 ภาษา ได้แก่ ภาษาไทยและภาษาอังกฤษเท่านั้น
    \item โมบายแอปพลิเคชันสามารถทำงานได้บนโทรศัพท์มือถือระบบปฏิบัติการแอนครอยด์ ตั้งแต่เวอร์ชัน \ldots{} เป็นต้นไป
    \item ด้วยข้อจำกัดด้านงบประมาณ ระบบจึงถูกติดตั้งและทดสอบบนเซิร์ฟเวอร์ของวิทยาลัยเท่านั้น\ldots{}
\end{enumerate}

% -----------------------------------------------------------
\section{ประโยชน์ที่คาดว่าจะได้รับ}
% -----------------------------------------------------------

\begin{enumerate}[label=\arabic*)]
    \item ได้ระบบ\ldots{} ที่สามารถ\ldots{}
    \item ได้แนวทางในการ\ldots{} ที่สามารถนำไปประยุกต์ใช้\ldots{}
    \item ได้องค์ความรู้ในด้าน\ldots{}
\end{enumerate}
